% فصل اول: مقدمه و کلیات پژوهش
\chapter{مقدمه و کلیات پژوهش}
\label{ch:intro}

بخش بزرگی از پدیده‌هایی که در علوم و مهندسی با آن‌ها سروکار داریم، از حرکت
سیالات و انتقال حرارت گرفته تا مکانیک کوانتومی و انتشار موج، با معادلات
دیفرانسیل جزئی\footnote{Partial Differential Equation (PDE)} توصیف می‌شوند.
به همین دلیل، حل دقیق و سریع این معادلات همواره یکی از مسائل بنیادی در
محاسبات علمی بوده است. از سوی دیگر، در سال‌های اخیر یادگیری ماشین و به‌ویژه
یادگیری عمیق\footnote{Deep Learning} موفقیت‌های چشمگیری در حوزه‌هایی مانند
بینایی ماشین و پردازش زبان به دست آورده است، اما به‌کارگیری مستقیم این
روش‌ها در مسائل فیزیکی معمولاً با چالش‌هایی مانند نیاز به داده آموزشی فراوان
و ناسازگاری پاسخ با قوانین فیزیک روبه‌روست. همین مسئله پژوهشگران را به این
فکر انداخته است که دانش فیزیکیِ مسئله، یعنی همان معادله دیفرانسیل حاکم بر
سیستم، را مستقیماً وارد فرایند یادگیری کنند. حاصل این ایده، خانواده‌ای از
روش‌ها با عنوان شبکه‌های عصبی آگاه از فیزیک\footnote{Physics-Informed Neural Network (PINN)} است که در
سال ۲۰۱۹ با مقاله رئیسی و همکارانش به‌طور جدی معرفی شد \cite{raissi2019} و
به‌سرعت به یکی از پراستنادترین موضوعات در مرز مشترک یادگیری ماشین و محاسبات
علمی تبدیل گردید.

در این پایان‌نامه، چارچوب شبکه‌های عصبی آگاه از فیزیک برای حل مسائل مستقیم و
معکوس شامل معادلات دیفرانسیل جزئی غیرخطی بررسی می‌شود. مسئله نمونه‌ای که در
بخش عملی این پژوهش انتخاب شده، معادله برگرز\footnote{Burgers Equation} است؛
معادله‌ای کلاسیک که به‌سبب رفتار غیرخطی و تشکیل جبهه‌های تند، معیاری شناخته‌شده
برای سنجش روش‌های عددی به شمار می‌رود. افزون بر مطالعه نظری، یک پیاده‌سازی
کامل از مدل زمان‌پیوسته با زبان پایتون و کتابخانه جکس\footnote{JAX} انجام و
نتایج آن با حل مرجع به دست آمده از یک حل‌گر طیفی مستقل مقایسه شده است.

\section{موضوع پژوهش}
\label{sec:subject}

موضوع این پایان‌نامه «شبکه‌های عصبی آگاه از فیزیک برای حل مسائل مستقیم و
معکوس معادلات دیفرانسیل جزئی غیرخطی» است. در مسئله مستقیم\footnote{Forward Problem}،
پارامترهای معادله معلوم فرض می‌شوند و هدف، یافتن جواب معادله با استفاده از
تعداد اندکی داده اولیه و مرزی است. در مسئله معکوس\footnote{Inverse Problem}،
برعکس، از روی مشاهده‌های پراکنده و احیاناً نویزیِ جواب، مقصود کشف پارامترهای
مجهول معادله حاکم بر سیستم است. هر دو دسته مسئله در مقاله مرجع این
پایان‌نامه \cite{raissi2019} بررسی شده‌اند و ما نیز هر دو را، هم از نظر نظری
و هم در قالب پیاده‌سازی عملی، دنبال می‌کنیم.

نکته مهم در این چارچوب، نقش مشتق‌گیری خودکار\footnote{Automatic Differentiation} است:
مشتق‌های جواب نسبت به متغیرهای مکان و زمان، به‌جای استفاده از شبکه محاسباتی
و تفاضلات محدود، مستقیماً از روی گراف محاسباتی شبکه عصبی و با دقت ماشین
محاسبه می‌شوند. به این ترتیب، باقیمانده معادله دیفرانسیل به‌صورت یک شبکه
عصبی جدید درمی‌آید که در همان تابع هزینه آموزش وارد می‌شود و شبکه را وامی‌دارد
پاسخی بیاموزد که هم به داده‌ها نزدیک باشد و هم معادله را ارضا کند.

\section{بیان مسئله}
\label{sec:problem}

روش‌های عددی کلاسیک مانند تفاضل‌های محدود\footnote{Finite Difference}،
المان‌های محدود\footnote{Finite Element} و روش‌های طیفی\footnote{Spectral Methods}
دهه‌هاست که ابزار اصلی حل معادلات دیفرانسیل جزئی هستند و مبانی نظری آن‌ها
به‌خوبی توسعه یافته است \cite{leveque2007}. با این حال، این روش‌ها معمولاً به
تولید شبکه محاسباتی وابسته‌اند، در ابعاد بالا با پدیده نفرین ابعاد\footnote{Curse of Dimensionality}
مواجه می‌شوند و برای هر هندسه یا شرط مرزی جدید باید از نو اجرا شوند. از طرفی،
مدل‌های یادگیری عمیقِ صرفاً داده‌محور هرچند در تقریب توابع بسیار توانمندند،
اما برای آموزش به حجم زیادی داده برچسب‌خورده نیاز دارند؛ داده‌ای که در
بسیاری از مسائل فیزیکی و مهندسی یا اصلاً در دسترس نیست یا تهیه آن بسیار
پرهزینه است. علاوه بر این، خروجی چنین مدل‌هایی هیچ تضمینی برای سازگاری با
قوانین فیزیک ندارد و ممکن است در خارج از محدوده داده آموزشی رفتار غیرفیزیکی
از خود نشان دهد.

مسئله‌ای که این پایان‌نامه به آن می‌پردازد را می‌توان چنین خلاصه کرد: چگونه
می‌توان دانش پیشینیِ موجود درباره یک سیستم فیزیکی، یعنی معادله دیفرانسیل
جزئی حاکم بر آن، را به‌صورت اصولی وارد آموزش یک شبکه عصبی کرد، به‌گونه‌ای که
اولاً با داده آموزشی اندک بتوان جواب دقیق مسئله را به دست آورد و ثانیاً از
روی مشاهده‌های محدود، پارامترهای مجهول خود معادله را نیز شناسایی کرد؟ پاسخ
مقاله مرجع به این پرسش، معرفی شبکه‌های عصبی آگاه از فیزیک در دو گونه مدل
زمان‌پیوسته و زمان‌گسسته است \cite{raissi2019} که مبنای نظری و عملی این
پژوهش را تشکیل می‌دهد.

\section{اهداف پژوهش}
\label{sec:goals}

هدف اصلی این پژوهش، مطالعه، پیاده‌سازی و ارزیابی چارچوب شبکه‌های عصبی آگاه
از فیزیک برای حل مسائل مستقیم و معکوس معادلات دیفرانسیل جزئی غیرخطی است.
برای دستیابی به این هدف اصلی، اهداف فرعی زیر دنبال می‌شوند:

\begin{enumerate}
\item مطالعه مبانی نظری لازم شامل معادلات دیفرانسیل جزئی، روش‌های عددی
کلاسیک، شبکه‌های عصبی، بهینه‌سازی و مشتق‌گیری خودکار؛
\item بررسی دقیق مقاله مرجع و استخراج مدل‌های زمان‌پیوسته و زمان‌گسسته برای
مسائل مستقیم و معکوس؛
\item پیاده‌سازی یک حل‌گر مرجع مستقل بر پایه روش طیفی فوریه برای معادله
برگرز به‌منظور تولید جواب دقیق جهت مقایسه؛
\item پیاده‌سازی مدل زمان‌پیوسته شبکه آگاه از فیزیک برای حل مستقیم معادله
برگرز و سنجش دقت آن نسبت به حل مرجع؛
\item پیاده‌سازی مسئله معکوس برای کشف ضرایب معادله برگرز از روی داده‌های
پراکنده و بررسی اثر نویز بر دقت شناسایی پارامترها.
\end{enumerate}

\section{سؤال‌های پژوهش}
\label{sec:questions}

سؤال‌هایی که این پژوهش در پی پاسخ‌گویی به آن‌هاست عبارت‌اند از:

\begin{enumerate}
\item آیا می‌توان با یک شبکه عصبی نسبتاً کوچک و تنها با استفاده از داده‌های
اولیه و مرزی، جواب دقیق یک معادله دیفرانسیل جزئی غیرخطی مانند معادله برگرز
را در کل دامنه زمانی ـ مکانی بازسازی کرد؟
\item سازوکار وارد کردن معادله دیفرانسیل در تابع هزینه شبکه چگونه است و
مشتق‌گیری خودکار در این میان چه نقشی ایفا می‌کند؟
\item دقت جواب به دست آمده از شبکه آگاه از فیزیک در مقایسه با یک حل‌گر عددی
استاندارد چقدر است و خطا در کدام نواحی دامنه بیشتر می‌شود؟
\item آیا می‌توان ضرایب مجهول معادله را نیز هم‌زمان با جواب، از روی
مشاهده‌های پراکنده و نویزی فراگرفت و دقت این شناسایی چه میزان است؟
\item مدل‌های زمان‌پیوسته و زمان‌گسسته هر یک برای چه دسته‌ای از مسائل
مناسب‌ترند و چه محدودیت‌هایی دارند؟
\end{enumerate}

\section{روش انجام پژوهش}
\label{sec:methodology}

این پژوهش از نظر روش، ترکیبی از مطالعه نظری و پیاده‌سازی تجربی است و در
مراحل زیر انجام شده است. نخست، مبانی نظری موضوع با مطالعه منابع معتبر در
زمینه معادلات دیفرانسیل جزئی \cite{evans2010}، روش‌های عددی
\cite{leveque2007} و یادگیری عمیق \cite{goodfellow2016} مرور شد. سپس مقاله
مرجع \cite{raissi2019} به‌طور کامل مطالعه و ایده‌های اصلی آن، یعنی تعریف
شبکه باقیمانده، طراحی تابع هزینه ترکیبی و دو گونه مدل زمان‌پیوسته و
زمان‌گسسته، استخراج گردید.

در مرحله بعد، معادله برگرز یک‌بعدی به‌عنوان مسئله نمونه انتخاب شد؛ زیرا هم
در مقاله مرجع به‌تفصیل بررسی شده و امکان مقایسه نتایج وجود دارد و هم به‌سبب
رفتار غیرخطی و تشکیل گرادیان‌های تند، مسئله‌ای چالشبرانگیز و در عین حال
قابل‌فهم برای سطح کارشناسی است. برای ارزیابی مستقل نتایج، یک حل‌گر مرجع با
روش طیفی فوریه و گام‌زنی زمانی نمایی مرتبه چهار پیاده‌سازی شد. آنگاه
مدل زمان‌پیوسته برای مسئله مستقیم و نیز مسئله معکوس کشف پارامترها با زبان
پایتون و کتابخانه جکس پیاده‌سازی و آموزش داده شد. در پایان، نتایج عددی شامل
دقت جواب، همگرایی آموزش و خطای شناسایی پارامترها تحلیل و با نتایج گزارش‌شده
در مقاله مرجع مقایسه گردید.

\section{ساختار پایان‌نامه}
\label{sec:structure}

ادامه این پایان‌نامه به شرح زیر سازماندهی شده است. در فصل دوم، مبانی نظری و
پیشینه پژوهش شامل معادلات دیفرانسیل جزئی، روش‌های عددی کلاسیک، شبکه‌های
عصبی، مشتق‌گیری خودکار و مروری بر کارهای پیشین در زمینه یادگیری ماشین آگاه
از فیزیک ارائه می‌شود. فصل سوم به تشریح روش پژوهش، یعنی چارچوب شبکه‌های
عصبی آگاه از فیزیک، مدل‌های زمان‌پیوسته و زمان‌گسسته و فرمول‌بندی مسائل
مستقیم و معکوس اختصاص دارد. در فصل چهارم، جزئیات پیاده‌سازی و نتایج عددی
برای معادله برگرز، شامل حل مستقیم و کشف پارامترها، گزارش و تحلیل می‌شود.
سرانجام در فصل پنجم، خلاصه‌ای از یافته‌ها، نتیجه‌گیری و پیشنهادهایی برای
ادامه پژوهش ارائه می‌گردد.
